\section{Open Challenges and Research Directions}
\label{sec:challenges}

\paragraph{1. Hearing below $-15$\,dB.}
On drones, even noise-model methods fail between $-15$ and $-20$\,dB (\cref{sec:case-dregon}), yet distant calls for help fall in exactly this range.
Smooth, physically grounded noise models that vary with each rotor's speed~\cite{yen2023ncm}, combined with learned suppression~\cite{wang2022drone,gulli2025rotor}, are the most promising route.
Hybrid methods that first remove tones predicted from the motor state and then whiten the remaining spatial noise exploit complementary structure.
Longer integration windows help when the target is persistent, but conflict with the motion of the platform; tracking with motion compensation (direction 5) can bridge the two.

\paragraph{2. From masking to cancellation with vibration references.}
Motor state gives the frequencies of rotor and motor tones but not their phase.
A high-rate vibration sensor near each motor, combined with order tracking, could supply a phase-coherent reference and allow tones to be subtracted rather than masked.
Standard flight-controller IMUs are sampled too slowly for this (\cref{sec:case-dregon}), so dedicated sensors are needed; to our knowledge this direction is largely unexplored on autonomous platforms.
Its value depends on the drone: the tonal share of ego-noise ranged from 1\% to 70\% across the five drones of \cref{sec:case-drones}. Rotor-speed variation in flight also smears the tones. Tonality should therefore be measured in flight before tone cancellation is pursued.

\paragraph{3. Transfer across platforms.}
Ego-noise models are platform-specific, and learned models trained on one drone or robot rarely transfer.
Our multi-drone analysis suggests a starting point: on every drone, a spatial noise model learned from the first half of a recording of tens of seconds remained valid on the second half (\cref{sec:case-drones}), so short calibration flights may suffice.
Self-supervised adaptation from such recordings, and models conditioned on platform geometry and state rather than identity, would reduce the cost of deploying audition on a new platform.
Array design matters as much as the algorithm: on a 6\,cm array a spatial model removed less ego-noise than it would remove from diffuse noise.
Large pretrained audio encoders~\cite{kong2020panns,gong2021ast,chen2023beats} are a natural starting point if they can be made robust to ego-noise.

\paragraph{4. Simulation with ego-noise.}
Embodied simulators render room acoustics well~\cite{chen2022soundspaces2,gao2023sonicverse,luo2022naf} but omit the agent's own noise.
Adding parametric ego-noise driven by the simulated platform's state, calibrated on real recordings such as DREGON or AIRA-UAS~\cite{strauss2018dregon,ruizespitia2018aira}, would make sim-to-real transfer of audio-visual policies far more credible.

\paragraph{5. Temporal reasoning, active listening and confidence.}
A platform should also decide \emph{which} method to trust: in our studies a noise model was essential on a drone but cost accuracy on a quiet robot head (\cref{sec:case-locata}), so selecting or weighting methods from the measured SNR and noise structure is a practical first step.
Most localizers decide each short window independently.
Tracking with the platform's own motion compensated~\cite{valin2007tracking,evers2018slam,evers2020locata} and calibrated confidence from the sharpness of the spatial spectrum would let an agent know when it is wrong.
An agent that knows it is uncertain can also act to listen better, by moving~\cite{nakadai2000active,majumder2021move2hear} or, on a drone, by briefly changing its rotor speeds.

\paragraph{6. Open-vocabulary hearing.}
Audio-language models~\cite{elizalde2023clap,deshmukh2023pengi,gong2024ltu,tang2024salmonn} could let a platform report what it hears in words and search for sounds described in language.
Their robustness to ego-noise, their latency and their energy cost on embedded hardware have yet to be measured.

\paragraph{7. Benchmarks that match the problem.}
A shared benchmark with a moving platform, real ego-noise, platform state, high-rate vibration sensing, realistic intermittent sources (calls, whistles, speech) and video (\cref{sec:datasets}), evaluated with the protocols of \cref{sec:metrics} and archived with a persistent identifier, would align robot audition, drone audition and embodied learning around one problem.

\paragraph{8. Efficiency and energy.}
Few drone localization systems run in real time on board~\cite{chevtchenko2025review}.
Method papers should report computational cost on embedded hardware alongside accuracy, as lightweight localization work already does~\cite{grondin2019lightweight}.

\paragraph{9. Privacy and responsible use.}
Always-on microphones on robots, drones and vehicles can record people who have not consented.
Processing on board, keeping only the outputs a task needs (events and directions rather than raw audio), and suppressing speech content when it is not needed are practical safeguards that systems and benchmarks should support.
